\section*{Use of AI Tools}
In the course of this academic work, the AI tools ChatGPT (GPT-4o, OpenAI) and DeepL Write were utilized to enhance translation, linguistic clarity, expression, and formatting. The tool was primarily used to enhance language and formatting, without generating or contributing to the academic content itself.

Specifically, AI assistance was applied during the writing of the following sections: \textit{Abstract}, \textit{Research (Presentation of Relevant Tools)}, \textit{Discussion} (\textit{Further Work}). In these sections, the tool was used exclusively to refine grammar, style, and readability.

Given that the thesis was written using LaTeX,  ChatGPT was also used to support formatting tasks, such as generating tables. The data used in these tables was originally created in Microsoft Excel and later formatted accordingly.

During the early research phase, ChatGPT was consulted to clarify basic concepts and terminology, supporting initial orientation and comprehension.

All substantive research, content development, analysis, and interpretation were carried out by the authors independently, based on personal work and literature review.

